\section{\vfast{}}
\label{sec:policy}

\subsection{Overview}

The FishBot v0.4 policy family is deterministic and fully inspectable, and has
four parts: an exact constraint tracker, a posterior over deals
(\S\ref{sec:belief}), a scored choice over legal asks, and a declaration
module. Each member observes only its own hand and the public record, and
there is no neural network, no tree search over sampled deals, and no private
channel between teammates: everything a teammate learns, an opponent learns at
the same instant.

The deployed configuration, evaluated throughout \S\ref{sec:results}, is
\vfast{}: its belief is the approximate path of \S\ref{sec:fast}, so it
consumes the per-card marginals $\hat\mu$ and acts on the allocation
probabilities $\hat\alpha$, while the exact reference engine of
\S\ref{sec:belief} runs only in validation and in the \vblock{} substitution
ablation, never in the deployed inner loop. Each equation below is labelled
fitted, approximate, or exact.

\subsection{Choosing an ask}
\label{sec:ask}

On its turn the policy scores every legal ask --- a pair $a = (c, t)$ naming a
card $c$ in a half-suit of which it holds at least one card, and an opposing
target $t$ --- and plays the maximiser of
\begin{equation}
  U(a) \;=\; \lambda_{\mathrm{lin}} \sum_{k=0}^{19} w_k\, \varphi_k(a)
       \;+\; \lambda_{\mathrm{val}}\,\Bigl[\, p\, V(s^{+}_a) + (1-p)\, V(s^{-}_a) \Bigr].
  \label{eq:askutil}
\end{equation}
Here $p = \hat\mu_{c,t}$ is the Fast posterior probability that the target
holds the named card; $\varphi(a) = (\varphi_0,\dots,\varphi_{19})$ is the
feature vector defined in Table~\ref{tab:features}; $w \in \mathbb{R}^{20}$
are the fitted ask weights; $\lambda_{\mathrm{lin}}$ and
$\lambda_{\mathrm{val}}$ are the fitted linear and value mixing weights; $V$
is the value function of \S\ref{sec:value}; and $s^{+}_a$, $s^{-}_a$ are the
positions reached after a hit and after a miss. The bracketed term is a
one-ply expectimax over $V$: a hit moves the card into our hand and retains
the turn, while a miss excludes the target --- renormalising that card's
ownership over the remaining candidates, so a miss is informative as well as
costly --- and hands the turn to the target.

Equation~\eqref{eq:askutil} is a \emph{fitted} score, not a derived quantity:
the weights $w$, $\lambda_{\mathrm{lin}}$ and $\lambda_{\mathrm{val}}$ come
from the search of \S\ref{sec:fitting}, and the expectation is taken against
$\hat\mu$ rather than the exact $\mu$. Ties break deterministically
(\S\ref{app:params}), and the chosen ask's $p$ is retained as the forecast
scored in the calibration study of \S\ref{sec:results-calib}.

The \numAskFeats{} coordinates of $\varphi$ are defined in
Table~\ref{tab:features}, placed in \S\ref{app:params} beside the weights
fitted to them so that a definition and its coefficient can be read together.
Five are representative. \emph{Hit probability} is $p$ itself, entered
linearly and squared so that the fit can curve the trade-off between certainty
and value; \emph{lock completion} is the probability that this ask alone
completes team ownership of the named card's half-suit $H$; \emph{reply
threat} is $(1-p)$ times the best ask the target could make against us if the
miss hands them the turn; and \emph{information leak}, with a companion
coordinate scaling it by how much of $H$ we already hold, indicates that our
team has never publicly revealed an interest in $H$. \emph{Runway} estimates
the length of the run this ask begins, since a hit retains the turn, nesting
the best remaining per-card posteriors as $p\,(1 + q_1(1 + q_2(1 + q_3)))$
with $q_1 \ge q_2 \ge q_3$ the next-best probabilities; it is an approximation
rather than an expectation over a modelled continuation, because it ignores
the opponents' replies and any change in the posterior along the run. All are
computed from the Fast marginals $\hat\mu$.

\subsection{Two-ply lookahead with approximate branch beliefs}
\label{sec:twoply}

Equation~\eqref{eq:askutil} evaluates both branches of an ask against the
\emph{current} posterior, and the continuation and runway features stand in
for what happens after it. For the leading candidates the policy instead
re-derives a belief on each branch: it ranks all legal asks by
\eqref{eq:askutil}, keeps the top $K$, and for each survivor builds the
post-hit and post-miss knowledge states (\S\ref{app:params}), propagates
capacities, and recomputes marginals on each.

The recomputation on both branches uses the same Sinkhorn/disjointness
approximation as the main path (\filepath{engine/src/v04.hpp}, \texttt{chooseAsk},
which calls \texttt{Belief::sinkhornDisj} on each branch state). Neither branch
belief is exact, and this operation is therefore not exact two-ply Bayesian
inference; it is a lookahead whose branch beliefs carry the same approximation
error as $\hat\mu$ itself, measured against the reference engine in
\S\ref{sec:results-verify}. Writing $\Pr'$ for the recomputed branch marginals,
the two summaries taken are
\begin{align*}
  \operatorname{follow}(a) &= \max_{a'} \Pr\nolimits'\bigl[a' \text{ hits}\bigr]
    \quad \text{over our legal asks in the post-hit position,} \\
  \operatorname{threat}(a) &= \max_{H} \Bigl( \Pr\nolimits'\bigl[t \text{ holds a card of } H\bigr]\cdot
        \max_{c' \in H} \Pr\nolimits'\bigl[c' \text{ is ours}\bigr] \Bigr)
    \quad \text{in the post-miss position,}
\end{align*}
and the candidate is rescored as
\begin{equation}
  U_2(a) \;=\; U(a) \;+\; \lambda_{\mathrm{chain}}\, p\,\operatorname{follow}(a)
       \;-\; \lambda_{\mathrm{threat}}\,(1-p)\,\operatorname{threat}(a),
  \label{eq:twoply}
\end{equation}
with the ask maximising $U_2$ played. The mixing weights
$\lambda_{\mathrm{chain}}$, $\lambda_{\mathrm{threat}}$ and the width $K$ are
all fitted coordinates of the parameter vector (\S\ref{sec:fitting}), so
\eqref{eq:twoply} is fitted and approximate on both counts. It is not a
determinized search: it stays inside the belief, for the reasons set out in
\S\ref{sec:related}.

In the frozen-policy ablations of \S\ref{sec:ablations}, removing the refinement
changed the win rate by \numAblTwoPly{} points with a 95\% paired interval of
\numAblTwoPlyCI{} (E5). The interval contains zero, so this experiment does not
resolve the refinement's contribution from zero.

\subsection{The value function}
\label{sec:value}

$V$ maps a vector of summary features of the current position to an estimate
of the final half-suit differential, scaled to $[-1,1]$ and signed from the
evaluating team's point of view. It is linear in \numValueFeats{} features ---
the score and card differentials, three control terms built from $e_H$, the
expected fraction of half-suit $H$ our team holds, counts of what is left to
play for, and side to move with its interactions --- defined and named in full
in Table~\ref{tab:valuecoef} of \S\ref{app:params}.

$V$ is fitted by ridge regression on self-play decision points labelled with
the eventual half-suit differential, with features collected from \emph{all
six seats} at every decision rather than from the mover alone; under
mover-only collection the side-to-move feature is constant at $+1$ and its
coefficient --- precisely what the miss branch of \eqref{eq:askutil} and the
stopping rule \eqref{eq:stop} depend on --- would be unidentified
(\S\ref{app:params}).

The fit reported as E14 has \numValueRows{} rows and an in-sample $R^2$ of
\numValueRSq{} (RMSE \numValueRmse{}) on the same self-play corpus that
generated the rows; no held-out evaluation of the value model was run. The
\numValueFeats{} coefficients compiled into the shipping policy are \emph{not}
the coefficients in that artifact, because the freeze step writes only the
\numParams{} policy parameters. The two vectors are printed side by side, with
the consequences of the gap, in \S\ref{app:params}; the metrics of E14
therefore describe the artifact's fit and not the coefficients that produced
the results in \S\ref{sec:results}.

\subsection{Declaring}
\label{sec:declare}

At every public event each agent is offered the chance to declare, so the full
posterior evaluation would run six times per event over every live half-suit.
Two cheap pre-gates screen candidates first: \texttt{Knowledge::cheapTeamProb}
scores a half-suit from capacities and supports alone, with no dynamic
program, against \texttt{gateTeamProb} ($0.008$), and if no live half-suit
clears it the policy declines to declare without building a posterior at all;
a second gate inside \texttt{evaluateSet} compares a product of per-card team
probabilities against \texttt{marginalGate} ($0.008$) and discards the
half-suit before the allocation query is run.

Neither gate is proved to preserve every candidate the ungated evaluation would
accept. They are pruning heuristics, and \S\ref{sec:results-verify} reports a
corpus-wide false-negative audit (E16) that measures the loss directly: rerunning
the full evaluation and the same stopping rule with both gates disabled found
\numGateFalseNeg{} false negatives among the rejected pairs, and the chosen
action differed at \numGateActionsPct\% of declaration opportunities. The gates
are therefore not lossless, and the audit bounds the size of the effect on the
primary evaluation corpus rather than eliminating it.

A half-suit that survives the gates is scored, and which estimator does the
scoring depends on the belief mode. \vfast{} computes $\hat\alpha$, the
probability that the named allocation is correct, by sequential conditioning
over the six cards (\S\ref{app:alpha}): the chaining is exact in structure,
and the approximation is entirely in the marginal solver. \vblock{} instead
reads the exact queries \eqref{eq:q2} and \eqref{eq:q3} straight off the block
tables, giving $\tau$ and $\alpha$. Under both modes the scored candidate is
then passed to the value-based stopping rule \eqref{eq:stop}.

Four quantities can override the stopping rule rather than consult it: the
policy is \emph{pressed} when the unresolved pool falls below
\texttt{patiencePool} cards, when the opposing team holds fewer than
\texttt{oppCardFloor} cards between them, when its own best available ask has
posterior below \texttt{askFloor}, or when the public event count passes the
forcing horizon --- the first three fitted coordinates of the parameter
vector, the horizon a fixed configuration constant. When pressed, a plain
probability threshold (the fitted \texttt{declThreshold}) replaces
\eqref{eq:stop}. The ask floor is a softened form of the dead-ask trigger that
\S\ref{sec:termination} reports as costly in its hard, unconditional form: as
a fitted floor combined with a threshold rather than a forced declaration, the
optimiser retained it at the value listed in Table~\ref{tab:params}. The rate
at which the floor fires was not measured.

Above all of this sits the graduated escalation of \S\ref{sec:termination},
which exists because two policies that both defer declaration were measured
failing to terminate within the action cap in mirror self-play
(Observation~\ref{obs:cycle}). Past its second horizon the pre-gates are
relaxed to match, so that a half-suit whose allocation is not sharply
estimated can still be named.

Two further decisions belong to the policy rather than to the driver: which
live teammate receives the turn when a declaration leaves a seat cardless, and
the willingness bit the policy supplies at each level of the forced-endgame
ladder of \S\ref{app:dialect-endgame}. Both are described in
\S\ref{app:params}.
